% Generálva: html/build_tex.js (ai11.js). Kézzel ne szerkeszd, a forrás a html/data mappa.
\setcounter{chapter}{10}
\chapter{\texorpdfstring{Forgalmi méretezés és VoIP}{Forgalmi méretezés és VoIP}}\label{f11}
Erlang-B táblázat vagy grafikon, exponenciális tartási idő, kihasználtság, és a VoIP csomagsebesség és sávszélesség számítása.

\section{Alapok, egyszerűen}
\subsection{\texorpdfstring{A1. Jelölések, Erlang-B}{A1. Jelölések, Erlang-B}}\label{a11-erlang}
\par\noindent{\small\renewcommand{\arraystretch}{1.25}\begin{tabularx}{\linewidth}{@{}>{\raggedright\arraybackslash}X>{\raggedright\arraybackslash}X@{}}\toprule
\textbf{Jel} & \textbf{Jelentés}\\\midrule
undefined & hívásintenzitás: átlagosan ennyi hívás érkezik időegységenként (Poisson)\\
undefined & átlagos tartási idő (exponenciális eloszlású)\\
undefined & felajánlott forgalom, undefined (Erlang)\\
undefined & kiszolgálók (vonalak, vevők, kocsik) száma\\
undefined & blokkolási valószínűség: a beérkező igény nem kap kiszolgálót\\
undefined & kihasználtság: egy kiszolgáló átlagos foglaltsága\\
\bottomrule\end{tabularx}}\par

\begin{center}\adjustbox{max width=\linewidth}{$\displaystyle P_b=\frac{A^N/N!}{\sum_{i=0}^NA^i/i!}\qquad a=\frac{A(1-P_b)}{N}$}\end{center}

Rekurzív alak (kézzel is gyors): $E_0=1$, $E_n=AE_{n-1}/(n+AE_{n-1})$, $P_b=E_N$. Táblázatnál a felajánlott forgalmat \textbf{felfelé} kerekítjük, és azt a legkisebb $N$-et keressük, amelynél $P_b$ még a határ alatt van. Exponenciális tartási időnél $P(X>t)=e^{-t/h}$.

\subsection{\texorpdfstring{A2. VoIP sávszélesség}{A2. VoIP sávszélesség}}\label{a11-voip}
$v_\text{forrás}$: a kodek bitsebessége, $l_p$: egy szegmens hossza (byte), $n$: szegmens csomagonként, $l_\text{fej}$: a fejrészek együtt (byte).

\begin{center}\adjustbox{max width=\linewidth}{$\displaystyle v_\text{cs}=\frac{v_\text{forrás}}{8\,n\,l_p}\qquad B=8\,(n\,l_p+l_\text{fej})\cdot v_\text{cs}$}\end{center}

Több szegmens egy csomagban: kevesebb fejrész, kisebb sávszélesség, de nagyobb csomagolási késleltetés, és egy elveszett csomag több beszédet visz el.

% ------------------------------------------------------------------
\section{\texorpdfstring{11.1. feladat: DTMF-vevők méretezése}{11.1. feladat: DTMF-vevők méretezése}}\label{ai11_01}
\setcounter{lepes}{0}
\begin{feladat}{11.1. feladat}
A tárcsázás átlagosan 10 s, óránként legalább 60 000 hívás, a blokkolás legfeljebb 0,2\%, exponenciális eloszlások. a) Felajánlott forgalom? b) Hány DTMF-vevő kell (a mellékelt táblázatból: 166 Erl $\to$ 197, 167 Erl $\to$ 198 a 0,2\%-os oszlopban)? c) Kihasználtság?

{\footnotesize\sffamily Forrás: ALL-IN 19., 39\ldots{}40. oldal $\cdot$ ai2016\_01\_05\_3, ai2016\_01\_05\_3b. Címkék: vizsga, javítás}
\end{feladat}
\begin{keresunk}
\begin{itemize}
\item a) $A$ (felajánlott forgalom)
\item b) $N$ (kiszolgálók száma)
\item c) $a$ (egy kiszolgáló kihasználtsága)
\end{itemize}
\end{keresunk}
\begin{jeloles}
\begin{jeltabla}
$\lambda$ & hívásintenzitás & \\
$\ca{h}$ & átlagos tartási idő & $10$\,s\\
$P_b$ & blokkolási valószínűség & \\
$A$ & felajánlott forgalom & \emph{keresett}\\
$N$ & kiszolgálók száma & \emph{keresett}\\
$a$ & egy kiszolgáló kihasználtsága & \emph{keresett}\\
\end{jeltabla}
\par\smallskip\textbf{Adott még:}
\begin{itemize}
\item $P_b\le0{,}2\%$ (blokkolási valószínűség)
\end{itemize}
\end{jeloles}

\lepes{a) Felajánlott forgalom}
$\lambda$-t és ${\ca{h}}$-t azonos időegységre kell hozni.

\begin{center}\adjustbox{max width=\linewidth}{$\displaystyle A=\lambda {\ca{h}}=\frac{60\,000}{3600\ \text{s}}\cdot10\ \text{s}=16{,}67\cdot10=166{,}67\ \text{Erl}$}\end{center}
\jegyzet{definíció}{}

\lepes{b) Vevők száma}
A táblázatban felfelé kerekítünk (167 Erl), különben a blokkolás a határ fölé kerülhet.

\begin{center}\adjustbox{max width=\linewidth}{$\displaystyle 167\ \text{Erl}\ \to\ N=198$}\end{center}
\jegyzet{számérték}{}

\begin{center}\adjustbox{max width=\linewidth}{$\displaystyle P_b(166{,}67;\ 197)=0{,}212\%>0{,}2\%\qquad P_b(166{,}67;\ 198)=0{,}178\%$}\end{center}
\jegyzet{ellenőrzés}{rekurzióval: 197 kevés}

\lepes{c) Kihasználtság}
Egy kiszolgáló kihasználtsága a lebonyolított forgalom, $A(1-P_b)$, osztva a kiszolgálók számával.

\begin{center}\adjustbox{max width=\linewidth}{$\displaystyle a=\frac{166{,}67\cdot(1-0{,}0018)}{198}=0{,}84$}\end{center}
\jegyzet{definíció}{}
\begin{figyelem}[Javítás]
Az egyik kézírásos változat 166-ra kerekít lefelé, és 197-et ad. A gépelt változat (167 $\to$ 198) helyes.
\end{figyelem}
\begin{eredmeny}

\begin{center}\adjustbox{max width=\linewidth}{$\displaystyle \text{a) }A=166{,}7\ \text{Erl}\qquad\text{b) }N=198\qquad\text{c) }a\approx84\%$}\end{center}
\end{eredmeny}


% ------------------------------------------------------------------
\section{\texorpdfstring{11.2. feladat: Autókölcsönző két kocsival}{11.2. feladat: Autókölcsönző két kocsival}}\label{ai11_02}
\setcounter{lepes}{0}
\begin{feladat}{11.2. feladat}
Naponta átlagosan 4 ügyfél (Poisson), átlagos bérlés 1,5 nap (exponenciális), 2 kocsi. a) Mekkora a valószínűsége, hogy nem tudnak kocsit adni? b) Átlagos kihasználtság? c) Napi haszon, ha egy kocsi fenntartása 60 USD/nap, a díj 100 USD/nap?

{\footnotesize\sffamily Forrás: ALL-IN 30\ldots{}31. oldal $\cdot$ ai2016\_06\_07\_7. Címkék: vizsga}
\end{feladat}
\begin{keresunk}
\begin{itemize}
\item a) $P_b$ (blokkolási valószínűség)
\item b) $a$ (egy kiszolgáló kihasználtsága)
\item c) napi haszon
\end{itemize}
\end{keresunk}
\begin{jeloles}
\begin{jeltabla}
$\lambda$ & hívásintenzitás & \\
$\ca{h}$ & átlagos tartási idő & $1{,}5$\\
$\cb{N}$ & kiszolgálók száma & $2$\\
$P_b$ & blokkolási valószínűség & \emph{keresett}\\
$a$ & egy kiszolgáló kihasználtsága & \emph{keresett}\\
$A$ & felajánlott forgalom & \\
\end{jeltabla}
\par\smallskip\textbf{Adott még:}
\begin{itemize}
\item fenntartás 60 USD/kocsi/nap, díj 100 USD/nap
\end{itemize}
\end{jeloles}

\lepes{a) Blokkolás}
Felajánlott forgalom, majd Erlang-B ${\cb{N}}=2$-re.

\begin{center}\adjustbox{max width=\linewidth}{$\displaystyle A=4\cdot{\ca{1{,}5}}=6\ \text{Erl}$}\end{center}
\jegyzet{definíció}{}

\begin{center}\adjustbox{max width=\linewidth}{$\displaystyle P_b=\frac{6^2/2!}{1+6+6^2/2!}=\frac{18}{25}=0{,}72$}\end{center}
\jegyzet{behelyettesítés}{}

\lepes{b) Kihasználtság}
Egy kocsi kihasználtsága a lebonyolított forgalom, $A(1-P_b)$ (ennyi kocsi van átlagosan kint), osztva a kocsik számával.

\begin{center}\adjustbox{max width=\linewidth}{$\displaystyle a=\frac{6\cdot(1-0{,}72)}{2}=\frac{1{,}68}{2}=0{,}84$}\end{center}
\jegyzet{definíció}{átlagosan 1,68 kocsi van kint}

\lepes{c) Haszon}
A bevétel a kint lévő kocsik átlagos száma szorozva a napi díjjal, a költség minden kocsira jár.

\begin{center}\adjustbox{max width=\linewidth}{$\displaystyle 1{,}68\cdot100-2\cdot60=168-120=48\ \text{USD/nap}$}\end{center}
\jegyzet{számérték}{}
\begin{eredmeny}

\begin{center}\adjustbox{max width=\linewidth}{$\displaystyle \text{a) }P_b=72\%\qquad\text{b) }a=84\%\qquad\text{c) }48\ \text{USD/nap}$}\end{center}
\end{eredmeny}
\begin{tipp}[Ellenőrzés]
\begin{itemize}
\item A gyűjtemény az a)-ra „passz”-t ír, pedig a b)-ben kiszámolja ugyanezt, a c) nincs megoldva.
\end{itemize}
\end{tipp}


% ------------------------------------------------------------------
\section{\texorpdfstring{11.3. feladat: 15 Erlang: hosszú hívások, kiszolgálók, VoIP}{11.3. feladat: 15 Erlang: hosszú hívások, kiszolgálók, VoIP}}\label{ai11_03}
\setcounter{lepes}{0}
\begin{feladat}{11.3. feladat}
Percenként 5 hívás, átlagosan 3 perc. a) Felajánlott forgalom? b) 100 hívásból hány hosszabb 6 percnél? c) Hány kiszolgáló kell 0,1\%-hoz? d) Kihasználtság? e) Elfér-e a forgalom VoIP-pal egy 10 Mbit/s-os vonalon?

{\footnotesize\sffamily Forrás: ALL-IN 46\ldots{}48., 85\ldots{}86. oldal $\cdot$ ai2016\_01\_12\_3, ai2016\_01\_12c\_3. Címkék: vizsga, javítás}
\end{feladat}
\begin{keresunk}
\begin{itemize}
\item a) $A$ (felajánlott forgalom)
\item b) $P(X>6\ \text{perc})$: annak valószínűsége, hogy egy hívás 6 percnél tovább tart
\item c) $N$ (kiszolgálók száma)
\item d) $a$ (egy kiszolgáló kihasználtsága)
\item e) VoIP-kapacitás
\end{itemize}
\end{keresunk}
\begin{jeloles}
\begin{jeltabla}
$\lambda$ & hívásintenzitás & \\
$\ca{h}$ & átlagos tartási idő & $3$\,perc\\
$P_b$ & blokkolási valószínűség & \\
$A$ & felajánlott forgalom & \emph{keresett}\\
$X$ & egy hívás tartási ideje (valószínűségi változó) & \emph{keresett}\\
$N$ & kiszolgálók száma & \emph{keresett}\\
$a$ & egy kiszolgáló kihasználtsága & \emph{keresett}\\
$E_0$ & a rekurzió kezdőértéke (0 kiszolgálóra), 1 & \\
$E_n$ & a blokkolási valószínűség $n$ kiszolgálóval & \\
$E_{n-1}$ & a blokkolási valószínűség $n-1$ kiszolgálóval & \\
$n$ & a kiszolgálók száma a rekurzió egy lépésében & \\
$E_N$ & a blokkolási valószínűség $N$ kiszolgálóval & \\
\end{jeltabla}
\par\smallskip\textbf{Adott még:}
\begin{itemize}
\item $P_b\le0{,}1\%$ (blokkolási valószínűség)
\end{itemize}
\end{jeloles}

\lepes{a) Forgalom}


\begin{center}\adjustbox{max width=\linewidth}{$\displaystyle A=5\cdot3=15\ \text{Erl}$}\end{center}
\jegyzet{definíció}{}

\lepes{b) Hosszú hívások}
Exponenciális eloszlásnál $P(X>t)=e^{-t/{\ca{h}}}$ (A1).

\begin{center}\adjustbox{max width=\linewidth}{$\displaystyle P(X>6)=e^{-6/3}=e^{-2}=0{,}135$}\end{center}
\jegyzet{behelyettesítés}{100-ból kb. 13 vagy 14}

\lepes{c) Kiszolgálók}
Erlang-B rekurzióval ($E_0=1$, $E_n=AE_{n-1}/(n+AE_{n-1})$, $P_b=E_N$) vagy táblázatból a legkisebb $N$, ahol $P_b\le0{,}1\%$.

\begin{center}\adjustbox{max width=\linewidth}{$\displaystyle N=27:\ P_b=0{,}160\%\qquad N=28:\ P_b=0{,}086\%$}\end{center}
\jegyzet{számérték}{$N=28$}

\lepes{d) Kihasználtság}
A lebonyolított forgalom, $A(1-P_b)$, osztva a kiszolgálók számával.

\begin{center}\adjustbox{max width=\linewidth}{$\displaystyle a=\frac{15\cdot0{,}999}{28}=0{,}54$}\end{center}
\jegyzet{definíció}{}

\lepes{e) VoIP}
Legfeljebb 28 egyidejű hívás. Egy G.711-es (64 kbit/s) hívás a fejrészekkel kb. 80 és 90 kbit/s között (A2).

\begin{center}\adjustbox{max width=\linewidth}{$\displaystyle 28\cdot87\ \text{kbit/s}\approx2{,}4\ \text{Mbit/s}<10\ \text{Mbit/s}$}\end{center}
\jegyzet{közelítés}{belefér, kétirányú forgalommal is}
\begin{figyelem}[Javítás]
A grafikonos változat 29 kiszolgálót olvas le, a közösségi gépelt kidolgozás a helyes 28-at.
\end{figyelem}
\begin{eredmeny}

\begin{center}\adjustbox{max width=\linewidth}{$\displaystyle \text{a) }15\ \text{Erl}\quad\text{b) }\approx14\quad\text{c) }N=28\quad\text{d) }a\approx54\%\quad\text{e) belefér}$}\end{center}
\end{eredmeny}


% ------------------------------------------------------------------
\section{\texorpdfstring{11.4. feladat: GSM-cella 1\%-os blokkolással}{11.4. feladat: GSM-cella 1-os blokkolással}}\label{ai11_04}
\setcounter{lepes}{0}
\begin{feladat}{11.4. feladat}
GSM, FRF = 12, a hívások átlagosan 2 percesek. Egy órában átlagosan hányan telefonálhatnak egy cellában, ha a blokkolás legfeljebb 1\%?

{\footnotesize\sffamily Forrás: ALL-IN 60\ldots{}62. oldal $\cdot$ ai2014\_12\_15\_9. Címkék: vizsga, javítás}
\end{feladat}
\begin{keresunk}
\begin{itemize}
\item $\lambda$: hívás/óra egy cellában
\end{itemize}
\end{keresunk}
\begin{jeloles}
\begin{jeltabla}
$\ca{h}$ & átlagos tartási idő & $2$\,perc\\
$P_b$ & blokkolási valószínűség & \\
$\lambda$ & hívásintenzitás & \emph{keresett}\\
$N$ & kiszolgálók száma & \\
$A$ & felajánlott forgalom & \\
\end{jeltabla}
\par\smallskip\textbf{Adott még:}
\begin{itemize}
\item $P_b\le1\%$ (blokkolási valószínűség)
\end{itemize}
\end{jeloles}

\lepes{Csatornák egy cellában}
Egy 12 cellás klaszter a GSM 900 mind a 124 vivőjét egyszer használja fel, vivőnként 8 időréssel, így egy cellára a klaszter csatornáinak tizenketted része jut.

\begin{center}\adjustbox{max width=\linewidth}{$\displaystyle \frac{124\cdot8}{12}=82{,}7\ \Rightarrow\ N=82$}\end{center}
\jegyzet{számérték}{a gyűjtemény 125 vivővel 83-at, kerekítve 80-at vesz}

\lepes{A megengedett forgalom}
Erlang-B: az a legnagyobb $A$, amelynél $P_b$ még 1\%.

\begin{center}\adjustbox{max width=\linewidth}{$\displaystyle N=82:\ A=67{,}2\ \text{Erl}\qquad(N=80:\ 65{,}4\qquad N=83:\ 68{,}2)$}\end{center}
\jegyzet{számérték}{}

\lepes{Hívásintenzitás}
$A=\lambda {\ca{h}}$-ból, ${\ca{h}}$ órában.

\begin{center}\adjustbox{max width=\linewidth}{$\displaystyle \lambda=\frac{A}{h}=\frac{67{,}2}{1/30\ \text{h}}=30\cdot67{,}2\approx2020\ \text{hívás/óra}$}\end{center}
\jegyzet{átalakítás}{}
\begin{figyelem}[Javítás]
A gyűjtemény 80 csatornával 65 Erl-t olvas le (ez jó), és 1950 hívás/órát kap. 83 csatornával 2045, 82-vel 2020 jön ki. Nagyságrendileg mind jó, a kerekítést mindig írd le.
\end{figyelem}
\begin{eredmeny}

\begin{center}\adjustbox{max width=\linewidth}{$\displaystyle \lambda\approx2000\ \text{hívás/óra (82 csatornával 2020)}$}\end{center}
\end{eredmeny}


% ------------------------------------------------------------------
\section{\texorpdfstring{11.5. feladat: Vállalati VoIP G.711-gyel (Moodle 2020)}{11.5. feladat: Vállalati VoIP G.711-gyel (Moodle 2020)}}\label{ai11_05}
\setcounter{lepes}{0}
\begin{feladat}{11.5. feladat}
10 ms-os szegmensekből 80 byte, csomagonként két szegmens, a fejrészek 47 byte. a) Csomagsebesség? b) A csatorna bitsebessége? c) Milyen hátránya lenne négy szegmensnek csomagonként?

{\footnotesize\sffamily Forrás: ALL-IN 14. oldal $\cdot$ ai2020\_06\_02\_Q6. Címkék: vizsga, javítás}
\end{feladat}
\begin{keresunk}
\begin{itemize}
\item a) $v_\text{cs}$ (csomagsebesség)
\item b) $B$ (egy hívás sávigénye)
\item c) $n=4$ (szegmensek száma csomagonként) hátránya
\end{itemize}
\end{keresunk}
\begin{jeloles}
\begin{jeltabla}
$\ca{l_p}$ & egy beszédszegmens hossza & $80$\,byte\\
$\cb{n}$ & szegmensek száma csomagonként & $2$\\
$l_\text{fej}$ & a fejrészek hossza & \\
$v_\text{cs}$ & csomagsebesség & \emph{keresett}\\
$B$ & egy hívás sávigénye & \emph{keresett}\\
$v_\text{forrás}$ & a kodek bitsebessége & \\
\end{jeltabla}
\end{jeloles}

\lepes{A forrás}


\begin{center}\adjustbox{max width=\linewidth}{$\displaystyle v_\text{forrás}=\frac{80\cdot8\ \text{bit}}{10\ \text{ms}}=64\ \text{kbit/s}$}\end{center}
\jegyzet{definíció}{G.711}

\lepes{a) Csomagsebesség}
Egy csomag két szegmens, azaz 20 ms beszéd.
\behely{\bh{\ca{l_p}}{\ca{80\un{byte}}}}

\begin{center}\adjustbox{max width=\linewidth}{$\displaystyle v_\text{cs}=\frac{64\,000}{8\cdot2\cdot{\ca{80}}}=\frac{1}{20\ \text{ms}}=50\ \text{csomag/s}$}\end{center}
\jegyzet{behelyettesítés}{}

\lepes{b) Bitsebesség}
Egy csomag $2\cdot{\ca{80}}+47=207$ byte.

\begin{center}\adjustbox{max width=\linewidth}{$\displaystyle B=207\cdot8\cdot50=82\,800\ \text{bit/s}=82{,}8\ \text{kbit/s}$}\end{center}
\jegyzet{behelyettesítés}{}

\lepes{c) Négy szegmens}
Kevesebb fej jut a beszédre, de a csomag csak 40 ms beszéd összegyűlése után indulhat, és egy csomag elvesztése 40 ms kiesést jelent.

\begin{center}\adjustbox{max width=\linewidth}{$\displaystyle 25\ \text{csomag/s},\ (320+47)\cdot8\cdot25=73{,}4\ \text{kbit/s}$}\end{center}
\jegyzet{számérték}{a sávnyereség kicsi, a késés kétszeres}
\begin{figyelem}[Javítás]
A hallgató 50 és 100 „bps”-t írt (mértékegység-tévesztés, a fejrészt is kihagyta). A javító 1 + 0 + 0 pontot adott.
\end{figyelem}
\begin{eredmeny}

\begin{center}\adjustbox{max width=\linewidth}{$\displaystyle \text{a) }50\ \text{csomag/s}\qquad\text{b) }82{,}8\ \text{kbit/s}\qquad\text{c) 40 ms csomagolási késleltetés}$}\end{center}
\end{eredmeny}


% ------------------------------------------------------------------
\section{\texorpdfstring{11.6. feladat: IP feletti média, fair queuing, WFQ}{11.6. feladat: IP feletti média, fair queuing, WFQ}}\label{ai11_06}
\setcounter{lepes}{0}
\begin{feladat}{11.6. feladat}
a) Mi nehezíti a „best effort” IP-n a multimédiát? b) Miért nem érdemes RTP helyett TCP-t használni? c) Mire való az IGMP? d) A fair queuing előnye a szigorú prioritással szemben, és a WFQ előnye?

{\footnotesize\sffamily Forrás: ALL-IN 28\ldots{}29., 74. oldal $\cdot$ ai2016\_06\_07\_5, ai2014\_01\_07\_7 (és 5-ösöm 29. oldal). Címkék: beugró}
\end{feladat}
\begin{keresunk}
\begin{itemize}
\item indokolt válaszok
\end{itemize}
\end{keresunk}
\begin{jeloles}
\par\smallskip\textbf{Adott még:}
\begin{itemize}
\item fogalmi kérdések
\end{itemize}
\end{jeloles}

\lepes{a) Best effort IP}
Változó késleltetés (jitter), csomagvesztés, sorrendcsere, és nincs garantált sávszélesség. A valós idejű média ezekre érzékeny.

\lepes{b) TCP vagy RTP}
A TCP az elveszett csomagot újraküldi: mire megérkezik, a lejátszásnál már késő. A torlódásvezérlése ráadásul visszafogja a sebességet. Az RTP (UDP felett) inkább eldobja a késő csomagot, és időbélyeget, sorszámot ad.

\lepes{c) IGMP}
Multicast csoporttagság kezelése: a vevő ezzel jelzi a routernek, melyik csoport (pl. IPTV-csatorna) forgalmát kéri.

\lepes{d) Ütemezés}
Szigorú prioritásnál a magas prioritású sor kiéheztetheti a többit. Fair queuing: folyamonként külön sor, körbejárva kiszolgálva, nincs kiéheztetés. WFQ: ugyanez súlyokkal, a folyamok a súlyuk arányában kapnak sávot, mindenki garantált minimummal.
\begin{eredmeny}

\begin{center}\adjustbox{max width=\linewidth}{$\displaystyle \text{a) jitter, vesztés}\quad\text{b) újraküldési késés}\quad\text{c) multicast}\quad\text{d) nincs kiéheztetés, WFQ: arányos részesedés}$}\end{center}
\end{eredmeny}
\begin{tipp}[Ellenőrzés]
\begin{itemize}
\item A gyűjtemény a)-ra a DF és MLR mérőszámokat írja: ezek az IPTV-minőség mérőszámai, nem az okok.
\end{itemize}
\end{tipp}
